\documentclass[12pt,a4paper]{article}

% ---- Encoding & Font ----
\usepackage{iftex}
\ifPDFTeX
  \usepackage[utf8]{inputenc}
  \usepackage[T1]{fontenc}
\else
  \usepackage{fontspec}  % XeTeX/LuaTeX (tectonic): Unicode nativo (·, —, ¿, ª, §)
\fi
\usepackage[spanish]{babel}
\usepackage{csquotes}

% ---- Page layout ----
\usepackage[top=2.5cm, bottom=2.5cm, left=2.5cm, right=2.5cm]{geometry}
\usepackage{setspace}
\onehalfspacing

% ---- Math ----
\usepackage{amsmath, amssymb, amsthm}
\usepackage{mathtools}
\usepackage{physics}
\usepackage{esint}  % \oiint

% ---- Graphics ----
\usepackage{graphicx}
\usepackage{tikz}
\usepackage{float}
\usepackage{caption}

% ---- Tables ----
\usepackage{array}
\usepackage{booktabs}
\usepackage{tabularx}
\usepackage{colortbl}

% ---- Colours ----
\usepackage{xcolor}
\definecolor{notebg}{HTML}{FFF3CD}
\definecolor{noteborder}{HTML}{FFC107}
\definecolor{boxred}{HTML}{F8D7DA}
\definecolor{boxredborder}{HTML}{F5C6CB}

% ---- Boxes ----
\usepackage{mdframed}
\usepackage{tcolorbox}
\tcbuselibrary{most}

\newtcolorbox{keybox}[1][]{
  colback=blue!5,
  colframe=blue!70!black,
  arc=4pt,
  boxrule=1pt,
  fonttitle=\bfseries,
  title={#1}
}

\newtcolorbox{warnbox}[1][]{
  colback=boxred,
  colframe=boxredborder,
  coltitle=black!75,
  arc=4pt,
  boxrule=1pt,
  fonttitle=\bfseries,
  title={#1}
}

\newtcolorbox{notebox}[1][]{
  colback=notebg,
  colframe=noteborder,
  coltitle=black!75,
  arc=4pt,
  boxrule=1pt,
  fonttitle=\bfseries,
  title={#1}
}

% ---- Hyperlinks & TOC ----
\usepackage[hidelinks]{hyperref}
\usepackage{bookmark}

% ---- Enumerate ----
\usepackage{enumitem}

% ---- Miscellaneous ----
\usepackage{icomma}
\usepackage{siunitx}
\usepackage{textcomp}
\usepackage[bottom]{footmisc}

% ---- Title ----
\title{\textbf{Clase 17: Campo de una Espira Circular, el Solenoide \\
        y la Ley de Ampère}\\[0.3em]
        \large{Campo sobre el eje, dipolo magnético, solenoide ideal
              y la ley de Ampère}}
\author{Transcripto y expandido --- Curso Física III (OpenFING)}
\date{}

\begin{document}

\maketitle
\thispagestyle{empty}
\newpage

\tableofcontents
\newpage

% ===================================================================
\section{Espira circular: campo sobre el eje}
% ===================================================================

Espira circular de radio $R$, corriente $I$; campo en $P$ sobre el eje a
distancia $Z$.

\subsection{Planteo y dirección del campo}

Por simetría, el elemento \textbf{diametralmente opuesto} tiene igual
componente axial pero componentes horizontales invertidas: al sumar,
\textbf{solo sobrevive la componente axial} $B_z$ (sentido por la mano
derecha).

\subsection{Cálculo de la componente axial}

Como $\mathbf{r}\perp d\mathbf{s}$:
\[
dB = \frac{\mu_0 I}{4\pi}\,\frac{ds}{r^2}, \qquad r^2 = Z^2 + R^2.
\]
Proyectando con $\sin\alpha = R/\sqrt{Z^2+R^2}$:
\[
dB_z = \frac{\mu_0 I R}{4\pi}\,\frac{ds}{(Z^2+R^2)^{3/2}}.
\]

\subsection{Resultado}

Nada depende del ángulo; $\int ds = 2\pi R$:

\begin{keybox}[Campo de una espira sobre su eje]
  \[
  \boxed{B_z = \frac{\mu_0\, I\, R^2}{2\,(Z^2+R^2)^{3/2}}}
  \]
\end{keybox}

\subsection{Casos límite: centro y dipolo magnético}

\textbf{Centro ($Z=0$):} $B_z = \mu_0 I/(2R) \neq 0$ (a diferencia del
anillo cargado). Las líneas son cerradas y atraviesan la espira.

\textbf{Lejos ($Z\gg R$):} $B_z \approx \dfrac{\mu_0 I R^2}{2|Z|^3} \sim
1/Z^3$.

\begin{keybox}[Dipolo magnético]
  Sin monopolos magnéticos, la espira vista de lejos es un \textbf{dipolo}
  ($1/r^3$, como el dipolo eléctrico), no una carga puntual ($1/r^2$). Se
  define el momento dipolar:
  \[
  \boxed{\mu = I\,A}, \qquad A = \pi R^2.
  \]
\end{keybox}

\begin{notebox}[Cuadrupolo]
  Si el dipolo se anula (dos espiras opuestas), sobrevive el efecto
  cuadrupolar ($1/r^4$).
\end{notebox}

% ===================================================================
\section{El solenoide}
% ===================================================================

\subsection{Definición y planteo}

Bobina cilíndrica con espiras \textbf{muy compactas}. Largo $L$, número
total $N$, densidad $n = N/L$ (espiras por unidad de longitud). Se calcula
$\mathbf{B}$ sobre el eje.

\begin{notebox}[Por qué compactas]
  Separadas serían hélices; pegadas, cada una se aproxima por una espira
  circular. El cable está barnizado (aislado) para que la corriente dé
  vueltas.
\end{notebox}

\subsection{Suma de espiras: la integral}

Una espira en $y$ (punto en $x$) aporta el resultado anterior con $Z \to
y-x$. En $dy$ hay $n\,dy$ espiras; integrando:
\[
B_z = \frac{\mu_0 I R^2 n}{2}\int_{-L/2}^{L/2}
      \frac{dy}{[(y-x)^2 + R^2]^{3/2}}.
\]
Con $u = y-x$ y la primitiva $\dfrac{u}{R^2\sqrt{u^2+R^2}}$:

\begin{keybox}[Campo de un solenoide sobre el eje]
  \[
  \boxed{B_z = \frac{\mu_0 n I}{2}\!\left[
    \frac{L/2 - x}{\sqrt{(L/2-x)^2 + R^2}} +
    \frac{L/2 + x}{\sqrt{(L/2+x)^2 + R^2}}\right]}
  \]
\end{keybox}

\subsection{El solenoide ideal}

$L \gg R$ con el punto \textbf{lejos de los bordes}: $|L/2 \mp x| \gg R$.
Tomando $L\to\infty$ con $x$ fijo, los cocientes tienden a $\pm 1$.

\subsection{Campo dentro y fuera}

\begin{keybox}[Solenoide ideal]
  \textbf{Dentro:} $\boxed{B_z = \mu_0\, n\, I}$ (constante a lo largo del
  eje). \textbf{Fuera (eje):} $B_z = 0$.
\end{keybox}

Líneas intensas y paralelas dentro; se abren al salir (alcance externo
$\sim R$).

\begin{notebox}[Solo sobre el eje]
  Este cálculo (Biot--Savart) es sobre el eje; la uniformidad en toda la
  sección se prueba con Ampère (\S3.4).
\end{notebox}

% ===================================================================
\section{Ley de Ampère}
% ===================================================================

\subsection{Enunciado}

Análogo magnético de Gauss (equivalencia con Biot--Savart no probada
aquí). Para una curva cerrada $C$:

\begin{keybox}[Ley de Ampère]
  \[
  \boxed{\oint_C \mathbf{B}\cdot d\mathbf{l} = \mu_0
         \sum_{k\ \text{a través de } C} I_k}
  \]
\end{keybox}

\begin{notebox}[Analogía con Gauss]
  Solo cuentan las corrientes que \textbf{atraviesan} $C$ (las exteriores
  generan campo pero no contribuyen a la circulación). Orientar las
  corrientes con la mano derecha según $C$. Útil solo con mucha simetría.
\end{notebox}

\subsection{Ejemplo: hilo infinito}

Simetría de rotación y traslación $\Rightarrow$ líneas circulares,
$\mathbf{B}$ tangente y de módulo constante. Con $C$ un círculo de radio
$R$:
\[
B\,(2\pi R) = \mu_0 I \;\Longrightarrow\; \boxed{B = \frac{\mu_0 I}{2\pi R}}.
\]
(Una sola línea de cálculo; para un segmento finito hay que usar
Biot--Savart.)

\subsection{Ejemplo: cable cilíndrico}

Cable de radio $R$, corriente $I$ uniforme. Con $\oint = B\,(2\pi r)$ y
la corriente encerrada:

\begin{keybox}[Cable cilíndrico]
  \[
  \boxed{B = \frac{\mu_0 I}{2\pi}\times
  \begin{cases}
  \dfrac{r}{R^2}, & r < R \\[2mm]
  \dfrac{1}{r}, & r > R
  \end{cases}}
  \]
\end{keybox}

Crece linealmente hasta $r=R$ (máximo $\mu_0 I/2\pi R$), luego decrece
como $1/r$.

\subsection{Ejemplo: solenoide ideal}

$\mathbf{B}$ paralelo al eje dentro, nulo fuera. Curva: rectángulo con un
lado (largo $a$) dentro y el opuesto lejos (fuera):
\begin{itemize}[nosep]
  \item Lados perpendiculares: $\mathbf{B}\perp d\mathbf{l}$, no aportan.
  \item Lado exterior: $B=0$.
  \item Lado interior: $\oint = B\,a$.
\end{itemize}
Corriente encerrada $I(n\,a)$:
\[
B\,a = \mu_0\, I\,(n\,a) \;\Longrightarrow\; \boxed{B = \mu_0\, n\, I}.
\]

\begin{keybox}[Vale en todo el interior]
  No se supuso estar sobre el eje: $B = \mu_0 n I$ vale en \textbf{cualquier
  punto interior} (lejos de los extremos), no solo sobre el eje.
\end{keybox}

\begin{notebox}[Próxima clase]
  Más aplicaciones de Ampère y materiales magnéticos.
\end{notebox}

\end{document}
